\documentclass[14pt,a4paper]{extarticle}

\usepackage{fontspec}
\defaultfontfeatures{Ligatures=TeX}
\setmainfont{Times New Roman}[
  Script=Cyrillic,
  Ligatures=TeX]
\newfontfamily\cyrillicfont{Times New Roman}[Script=Cyrillic]

\usepackage{polyglossia}
\setdefaultlanguage{russian}
\usepackage{geometry}
\geometry{top=18mm,bottom=18mm,left=20mm,right=15mm}
\usepackage{setspace}
\usepackage{microtype}
\usepackage{indentfirst}
\usepackage{hyperref}
\hypersetup{hidelinks}

\setstretch{1.0}
\pagestyle{empty}
\setlength{\parindent}{0pt}
\setlength{\parskip}{0pt}
\sloppy

\newcommand{\annotationsection}[2]{%
  \textbf{#1 / #2}\par
}

\begin{document}

\annotationsection{Цель исследования}{Research goal}
Разработать микросервис, поддерживающий актуальную таблицу соответствия
IP-префиксов автономным системам на основе данных протокола BGP в режиме
реального времени для потребителей DDoS-защиты и сетевой аналитики.

\annotationsection{Задачи, решаемые в ВКР}{Research tasks}
\begin{enumerate}
\item изучить предметную область междоменной маршрутизации, свойства BGP
и проблему свежести маршрутного состояния;
\item проанализировать известные подходы к LPM и структурам хранения
префиксных таблиц;
\item сформулировать функциональные и нефункциональные требования к
сервису, работающему по схеме «начальный снимок + поток обновлений»;
\item спроектировать архитектуру микросервиса, реализовать gRPC API для
потребителей и интерфейсы взаимодействия с узлом GoBGP как источником
маршрутных данных;
\item реализовать загрузку начального снимка и потоковую обработку
обновлений с восстановлением после разрыва соединения;
\item реализовать шардированное in-memory хранилище на базе
Adaptive Radix Tree с операциями поиска, вставки, удаления и полной
выгрузки;
\item провести сравнительный эксперимент с прежним вариантом на ranger
и оценить компромиссы между временем чтения, скоростью обновления и
временем полной выгрузки.
\end{enumerate}

\annotationsection{Краткая характеристика полученных результатов}{Short summary of results/findings}
В рамках работы был разработан микросервис, поддерживающий таблицу соответствия IP-префиксов автономным системам в 
актуальном состоянии на основе данных BGP и предоставляющий единый gRPC API для нескольких внутренних потребителей.
Реализована модель «начальный снимок + поток обновлений»: сервис
получает исходное состояние от GoBGP, применяет одиночные BGP-события к
общей in-memory базе и выполняет пересинхронизацию после разрыва
соединения. Логика построения и обновления таблицы вынесена из состава
BIFIT COLLECTOR в отдельный сервис, пригодный для совместного
использования BIFIT COLLECTOR и BIFIT MITIGATOR. Для хранения состояния
реализовано шардированное хранилище на основе Adaptive Radix Tree,
поддерживающее операции Lookup, Upsert, Delete, Replace и Dump.
Экспериментальное сравнение с прежним вариантом на ranger показало, что
новое хранилище существенно ускоряет одиночные инкрементальные
обновления и делает потоковое применение BGP-событий эксплуатационно
осуществимым. При этом ranger сохраняет преимущество для редких
изменений состояния с частыми операциями Lookup и Dump, что задает
границу применимости полученного решения.

\end{document}
